\chapter{Chi viveva sulle Alpi, e come si comportava}
\label{cap:specie_alpine}

Un cacciatore preistorico non sceglieva dove accamparsi in astratto: sceglieva in funzione di dove si trovava la sua preda, e di come quella preda si muoveva. Per capire dove cercare le tracce dei cacciatori, allora, bisogna prima capire chi cacciavano, e soprattutto quanto e come si spostavano gli animali che cacciavano. Questo capitolo raccoglie i dati quantitativi disponibili sul comportamento delle specie alpine più rilevanti, con le fonti che li sostengono, perché è proprio la precisione di questi numeri, non solo la loro esistenza, a permettere le scelte metodologiche della terza parte di questo libro.

\section{Il camoscio: un territorio piccolo e fedele}

Il camoscio (\textit{Rupicapra rupicapra}) è la specie meglio studiata tra gli ungulati alpini, e anche quella con il territorio più ridotto. Una revisione sistematica recente, che ha raccolto ventidue studi condotti in sedici aree di studio diverse, riporta home range annuali individuali compresi tra 0,04 e 4,94 chilometri quadrati, con una variabilità enorme legata al sesso (i maschi occupano territori più ampi delle femmine), al comportamento di dispersione (più ampio nei maschi migratori) e alla stagione \citep{papakostas2026}. Le femmine, in particolare, mostrano territori particolarmente ristretti: tra 0,64 e 0,87 chilometri quadrati, con una media di 0,72.

Uno studio condotto nel Parco Nazionale del Gran Paradiso, sulle Alpi Graie, offre un dettaglio ancora più preciso e utile per il ragionamento di questo libro: distingue tra maschi residenti (con home range mediano di quarantanove ettari) e maschi migratori (home range mediano di settecentoquarantanove ettari, quindici volte più grande), e mostra che, nella stagione calda, le dimensioni del territorio aumentano significativamente rispetto alla stagione fredda per entrambe le categorie. Un altro studio, condotto nelle Alpi bavaresi, riporta valori medi ancora diversi: tra centoquarantaquattro e seicentotrentacinque ettari, con una media di trecentonovantotto.

Questa variabilità tra studi non è un difetto dei dati, ma un'informazione in sé: mostra che il territorio del camoscio dipende fortemente dal contesto locale (qualità dell'habitat, presenza di predatori, densità di popolazione), e che qualunque cifra usata in questo libro va trattata come un ordine di grandezza, non come un valore universale. Ma anche tenendo conto di questa variabilità, un punto resta fermo in tutti gli studi consultati: il camoscio occupa un territorio piccolo, quasi sempre inferiore ai cinque chilometri quadrati, spesso molto meno.

\section{Lo stambecco: territorio più ampio, ma fedele nel tempo}

Lo stambecco (\textit{Capra ibex}) occupa territori più estesi del camoscio, ma condivide con esso un tratto comportamentale importante: la fedeltà al proprio spazio nel corso degli anni. La stessa revisione sistematica citata sopra riporta, per lo stambecco, un intervallo di home range annuale tra 6,1 e 22,6 chilometri quadrati.

Uno studio di lungo periodo condotto ancora nel Gran Paradiso mostra un pattern di uso dello spazio fortemente tradizionale: gli individui tornano, anno dopo anno, quasi esattamente negli stessi luoghi, con una sovrapposizione elevatissima tra il territorio usato in anni consecutivi. La copertura nevosa limita fortemente i movimenti invernali, tanto che gli home range invernali e primaverili risultano sistematicamente più piccoli di quelli usati tra estate e autunno.

Questo comportamento è particolarmente rilevante per il tema di questo libro: uno stambecco non è un animale che percorre lunghe distanze in cerca di nuovi territori, ma un animale che ritorna, quasi con precisione geografica, sempre nelle stesse zone. Se un gruppo di cacciatori avesse individuato, in una data stagione, la posizione di una colonia di stambecchi, avrebbe avuto buone ragioni per aspettarsi di trovarla, l'anno successivo, quasi nello stesso luogo.

\section{Il cervo: l'eccezione migratrice}

Il cervo (\textit{Cervus elaphus}) si comporta in modo sostanzialmente diverso dalle due specie precedenti, ed è l'unica, tra quelle qui considerate, per cui ha senso parlare di una vera migrazione stagionale su distanze significative.

Due studi condotti nelle Alpi bavaresi da Bernhard Georgii, uno sulle femmine e uno sui maschi, offrono probabilmente il quadro più dettagliato disponibile. Per le femmine, l'home range medio misura sessantacinque ettari in inverno, centosessantasette ettari tra primavera e autunno, e centoventuno ettari in estate \citep{georgii1980a}: non tutte le femmine migrano allo stesso modo, alcune restano tutto l'anno in pianura, mentre altre si spostano stagionalmente tra due o anche tre aree diverse. Per i maschi, l'home range invernale e quello del periodo degli amori (l'accoppiamento autunnale) sono relativamente piccoli, in media rispettivamente centotredici e centotrentaquattro ettari, mentre l'home range usato tra primavera e autunno sale fino a circa trecentottantasei ettari \citep{georgii1981}.

Uno studio più recente condotto nelle Alpi italiane conferma questo quadro su scala diversa, con home range annuali complessivi tra 0,2 e 68,7 chilometri quadrati a seconda dell'individuo, e con una differenza di comportamento tra popolazioni prevalentemente sedentarie e popolazioni migratorie che occupano quote basse in inverno e quote più elevate in estate \citep{luccarini2006}.

\section{Uno sguardo d'insieme}

La tabella seguente raccoglie i valori discussi in questo capitolo, per permettere un confronto diretto tra le tre specie.

\begin{table}[htbp]
\centering
\small
\begin{tabular}{@{}p{2.6cm}p{3.6cm}p{2.2cm}p{2.6cm}@{}}
\toprule
\textbf{Specie} & \textbf{Home range annuale} & \textbf{Variazione stagionale} & \textbf{Fonte} \\
\midrule
Camoscio & 0,04--4,94 km$^2$ (femmine: 0,64--0,87 km$^2$) & sì, minore in inverno & Papakostas et al., 2026 \\
Stambecco & 6,1--22,6 km$^2$ & sì, minore in inverno, forte fedeltà al sito & Papakostas et al., 2026 \\
Cervo (femmine) & 0,65--1,67 km$^2$ (65--167 ha) & sì, marcata & Georgii, 1980 \\
Cervo (maschi) & 1,1--3,9 km$^2$ (113--386 ha) & sì, marcata & Georgii, 1981 \\
\bottomrule
\end{tabular}
\caption{Home range delle tre principali specie di ungulato alpino discusse in questo capitolo.}
\label{tab:home_range}
\end{table}

Il quadro che emerge, se si mettono insieme questi dati, è netto: camoscio e stambecco occupano territori piccoli e li riusano fedelmente nel tempo, con un restringimento invernale legato alla neve; il cervo occupa territori generalmente più ampi e, in una parte della popolazione, compie vere migrazioni stagionali tra quote diverse. Questa distinzione, tra specie territorialmente fedeli e a raggio ridotto e una specie più mobile con vera migrazione stagionale, tornerà al centro della discussione nella terza parte di questo libro, quando si affronterà la domanda se sia sensato immaginare, per i cacciatori preistorici, dei corridoi di spostamento lunghi e continui, oppure territori di caccia più piccoli e stabili, riutilizzati stagione dopo stagione.

\section{Una fauna che cambia nel tempo}

Un secondo punto, altrettanto importante quanto i dati di home range appena discussi, è che la composizione della fauna alpina non è rimasta costante per tutto il periodo che interessa questo libro. È cambiata, e in modo significativo, seguendo proprio le oscillazioni climatiche ricostruite nel capitolo precedente.

Nelle fasi più antiche, quando il clima era ancora più freddo e il paesaggio più aperto, tra le prede compaiono anche specie oggi associate ad ambienti di steppa o di prateria fredda: l'uro, un grande bovino selvatico oggi estinto, progenitore del bestiame domestico; il bisonte delle steppe; il cavallo selvatico. Sono animali di ambiente aperto, che frequentano fondovalle ampi e pianure, non i valichi stretti e le creste che si associano più naturalmente a stambecco e camoscio, come mostrerà nel dettaglio, con i dati archeozoologici dei siti prealpini veneti, il prossimo capitolo.

Con il progressivo miglioramento climatico e l'avanzata delle foreste descritta nel capitolo precedente, questo quadro cambia. Le grandi specie di ambiente aperto scompaiono progressivamente dal registro archeologico locale, sostituite da una fauna più simile a quella discussa in questo capitolo: stambecco e camoscio diventano le prede dominanti, affiancate da cervo e capriolo; compare, in questa fase più recente, anche l'alce, specie legata ad ambienti boschivi e umidi che si diffonde solo dopo l'insediamento definitivo delle foreste.

Questo significa che un modello di ricerca pensato per un momento preciso di questa lunga vicenda rischia di non funzionare per un momento diverso. Un sito di caccia all'uro, animale di pianura e di ambiente aperto, va cercato in un luogo molto diverso da un sito di caccia allo stambecco, animale di quota e di terreno impervio. Qualunque criterio di ricerca deve dichiarare, fin da subito, a quale fase e a quale tipo di fauna si riferisce: non esiste un criterio unico buono per tutto l'arco cronologico che va dalle prime frequentazioni fino alla comparsa dell'agricoltura.

\section{Un avvertimento sui dati moderni}

I valori di home range riportati nella tabella~\ref{tab:home_range} provengono da studi condotti su popolazioni di ungulati contemporanee, monitorati con radiotelemetria e GPS in ambienti alpini attuali. Prima di usarli per ragionare sul comportamento delle stesse specie diecimila anni fa, è necessario dichiarare esplicitamente un limite che questi studi non possono colmare.

La pressione di predazione che subivano camoscio e stambecco nel Mesolitico era radicalmente diversa da quella attuale. Il lupo era abbondante sull'intero arco alpino; la lince era presente; l'orso bruno, pur essendo principalmente onnivoro, può predare ungulati. E a queste pressioni naturali si aggiungeva quella dei cacciatori mesolitici specializzati, che proprio stambecco e camoscio avevano come preda principale, come mostra il capitolo~\ref{cap:cosa_cacciavano}. Una popolazione di ungulati cacciata sistematicamente da predatori naturali e da cacciatori umani, stagione dopo stagione sullo stesso territorio, può sviluppare pattern di movimento molto diversi da quelli osservati in una popolazione moderna in cui il lupo è stato assente per secoli e la caccia sportiva è regolamentata e stagionale.

Questo non significa che i valori di home range moderni siano inutili. Significa che vanno usati come ordini di grandezza e come indici di comportamento qualitativo (territorio piccolo e fedele nel tempo vs. territorio ampio e migratorio), non come cifre precise trasferibili senza correzione a un contesto di diecimila anni fa. Dove in questo libro si usano questi dati per costruire un'ipotesi su dove si trovavano le prede dei cacciatori mesolitici, quella conclusione è un'analogia ragionevole, non una certezza dimostrata.
